% Source: https://artofproblemsolving.com/community/c3h1767962_whats_the_correct_answer_of_this_question
% Stable identity: aops-topic-1767962-q01
Let $g(t)=\sqrt{3+t}$ for $t\ge-3$; the equation is $g(g(g(x)))=x$. The left side is a square root, so $x\ge0$, and $g$ is strictly increasing.

If $g(x)>x$, then applying the increasing function $g$ gives $g(g(x))>g(x)>x$ and then $g(g(g(x)))>g(g(x))>x$, a contradiction. Likewise $g(x)<x$ leads to $g(g(g(x)))<x$. Hence every solution satisfies $g(x)=x$, i.e.\ $\sqrt{3+x}=x$, so $x^2-x-3=0$ with $x\ge0$, giving $x=\dfrac{1+\sqrt{13}}{2}$. Conversely this value satisfies $g(x)=x$ and hence the equation.

There is exactly $\boxed{1}$ real solution, $x=\dfrac{1+\sqrt{13}}{2}\approx2.30$.
